\documentclass[14pt]{ca-sheet}
% Generated by taskgen from sheet "test-legacy" (master seed 4242).
% Do not edit: regenerate with `taskgen regen <manifest.json>`.
\begin{document}
\sheettitle{Тест}


\clearpage

\tableofcontents

\clearpage
\section{Варианты группы A 1}
% variant A 1/1
\subsection{Вариант 1}
\begin{enumerate}
  \item Решить задачу дискретного логарифмирования методом Шенкса: \[29^x \equiv 26\ (\textrm{mod}\ 83).\] Обосновать выбор длины списков.
  \item Вычислить степень $19P$ точки $P = (52, 28)$ эллиптической кривой \[y^2 \equiv x^3 + 88x + 33\ (\textrm{mod}\ 97)\] с помощью тернарного разложения показателя $19$.
\end{enumerate}
\clearpage
% variant A 1/2
\subsection{Вариант 2}
\begin{enumerate}
  \item Возвести целое число в степень по модулю с помощью бинарного алгоритма: \[115^{33}\ (\textrm{mod}\ 336).\]
  \item Вычислить степень $27P$ точки $P = (9, 54)$ эллиптической кривой \[y^2 \equiv x^3 + 44x + 11\ (\textrm{mod}\ 89)\] с помощью бинарного разложения показателя $27$.
\end{enumerate}
\clearpage
\section{Варианты группы B}
% variant B/1
\subsection{Вариант 1 --- Иванов}
\begin{enumerate}
  \item Решить задачу дискретного логарифмирования методом Шенкса: \[25^x \equiv 56\ (\textrm{mod}\ 113).\] Обосновать выбор длины списков.
  \item Вычислить степень $11P$ точки $P = (68, 14)$ эллиптической кривой \[y^2 \equiv x^3 + 23x + 8\ (\textrm{mod}\ 71)\] с помощью бинарного разложения показателя $11$.
\end{enumerate}
\clearpage
% variant B/2
\subsection{Вариант 2 --- Петров}
\begin{enumerate}
  \item Возвести целое число в степень по модулю с помощью бинарного алгоритма: \[911^{22}\ (\textrm{mod}\ 424).\]
  \item Вычислить степень $47P$ точки $P = (64, 76)$ эллиптической кривой \[y^2 \equiv x^3 + 42x + 31\ (\textrm{mod}\ 97)\] с помощью тернарного разложения показателя $47$.
\end{enumerate}
\clearpage
\end{document}
